\subsection{Diseño experimental}

\subsubsection{Estrategias comparadas}

El experimento se organizó alrededor de tres estrategias de representación del vector de estado: una referencia sin compresión, la propuesta sin pérdida basada en Blosc y una alternativa con pérdida controlada basada en ZFP. Las tres se ejecutaron sobre los mismos circuitos, con iguales tamaños de qubits y bajo un procedimiento homogéneo de medición, de modo que las diferencias observadas pudieran atribuirse al esquema de almacenamiento y no a cambios en la carga, en la instrumentación o en el entorno de ejecución.

\subsubsection{Selección de cargas de trabajo}

La evaluación se estructuró sobre dos familias de circuitos con comportamientos contrastantes en términos de compresibilidad: el algoritmo de búsqueda de Grover y la transformada cuántica de Fourier (QFT). Esta selección evitó sesgar la evaluación hacia un único tipo de estado cuántico. Grover representa escenarios en los que pueden mantenerse regularidades aprovechables por la compresión, mientras que QFT permite examinar casos en los que esa estructura se reduce o desaparece. El diseño, por tanto, cubre tanto condiciones favorables para la propuesta como situaciones menos propicias, necesarias para delimitar su alcance real.

\paragraph{Casos evaluados para Grover}

En Grover se estudiaron dos variantes. La primera correspondió a un oráculo con objetivo único y se ejecutó con tres ubicaciones representativas del estado marcado: $N/8$, $N/2$ y $7N/8$, donde $N = 2^n$ es el tamaño del espacio de búsqueda. La segunda correspondió a un oráculo multiobjetivo con tres estados marcados distribuidos en esas mismas regiones. Con esta definición se buscó que los resultados no dependieran de una sola disposición de índices y, al mismo tiempo, observar si la posición de los elementos marcados introducía variaciones apreciables en memoria o tiempo de ejecución.

\paragraph{Casos evaluados para QFT}

En QFT se consideraron dos familias de entrada. La primera fue una superposición periódica, incluida por su estructura regular y repetitiva. La segunda fue una superposición de alta entropía con amplitudes de magnitud uniforme y fases aleatorias, definida como $a_y = e^{i\phi_y}/\sqrt{N}$ con $\phi_y \sim U[0, 2\pi)$ y semilla fija. Este contraste permitió enfrentar la propuesta tanto a estados con redundancia potencialmente explotable como a estados donde dicha redundancia es escasa, mientras que el uso de una semilla fija preservó la reproducibilidad del experimento.

\subsubsection{Escala del experimento}

La evaluación se ejecutó para tamaños entre 20 y 25 qubits. Este rango permitió trabajar con instancias suficientemente grandes para que la memoria se convirtiera en una restricción relevante dentro de la capacidad disponible del entorno experimental, sin perder reproducibilidad ni comprometer la comparabilidad entre estrategias. Elegir varios tamaños dentro de ese intervalo también permitió observar cómo evolucionaban los efectos de la compresión a medida que crecía el vector de estado.

\subsubsection{Entorno experimental e instrumentación}

Todas las ejecuciones se realizaron en una plataforma experimental con Rocky Linux 10.1, procesador AMD Ryzen 7 5700X de 8 núcleos y 16 hilos, y 32~GiB de memoria RAM. La recolección de métricas se efectuó con los scripts de perfilado del repositorio del simulador, que ejecutan los binarios experimentales y registran tiempo transcurrido y memoria máxima mediante \texttt{psrecord}. Se recurrió a \texttt{psrecord} porque, al operar como una herramienta externa al código instrumentado, permite observar el consumo de recursos del proceso sin introducir lógica adicional en la implementación evaluada. Con ello se buscó reducir la interferencia de la propia medición sobre el experimento y mantener un procedimiento uniforme para todas las estrategias comparadas.

\subsubsection{Configuración de compresión}

Las estrategias comprimidas se evaluaron con configuraciones fijas durante toda la campaña experimental, resumidas en la Tabla~\ref{tab:exp-configs}. Esta decisión se adoptó después de una fase preliminar de ajuste, separada de la evaluación principal, en la que se exploraron valores de granularidad de bloque, tamaño de caché y parámetros de compresión sobre ejecuciones representativas. El criterio de selección fue privilegiar configuraciones que mantuvieran ahorro de memoria positivo cuando existía estructura compresible y que, simultáneamente, evitaran una sobrecarga temporal excesiva frente a la referencia. Una vez fijadas, las configuraciones se conservaron sin cambios para impedir que la comparación quedara condicionada por un reajuste particular en cada caso.

En particular, los valores $\texttt{chunkStates}=32768$ y $\texttt{gateCacheSlots}=8$ se eligieron porque ofrecieron una granularidad suficiente para amortiguar el costo del códec sin forzar materializaciones demasiado grandes, y porque permitieron retener en memoria de trabajo los bloques más reutilizados durante la aplicación de compuertas. En Blosc, $\texttt{clevel}=1$ y $\texttt{useShuffle}=\texttt{true}$ se fijaron por representar una configuración conservadora, favorable al tiempo de ejecución y coherente con arreglos homogéneos de dobles intercalados. En ZFP, el modo \texttt{FixedPrecision} con \texttt{precision}=40 se mantuvo como punto de comparación con pérdida controlada, priorizando una discrepancia numérica pequeña sin eliminar por completo el potencial de reducción de memoria.

\begin{table}[H]
    \centering
    \small
    \caption{Configuraciones de compresión utilizadas en toda la campaña experimental}
    \label{tab:exp-configs}
    \begin{tabular}{>{\raggedright\arraybackslash}p{2.1cm}>{\raggedright\arraybackslash}p{9.2cm}}
        \hline
        \textbf{Estrategia} & \textbf{Configuración} \\
        \hline
        Blosc & $\texttt{chunkStates}=32768$, $\texttt{gateCacheSlots}=8$, $\texttt{clevel}=1$, $\texttt{nthreads}=4$, $\texttt{compcode}=1$, $\texttt{useShuffle}=\texttt{true}$. \\
        ZFP & Modo $\texttt{FixedPrecision}$, $\texttt{precision}=40$, $\texttt{chunkStates}=32768$, $\texttt{gateCacheSlots}=8$, $\texttt{nthreads}=4$. \\
        \hline
    \end{tabular}
\end{table}

\subsubsection{Métricas de evaluación}

Las métricas principales del estudio fueron el tiempo total de ejecución y la memoria máxima. Ambas se obtuvieron a partir de las trazas generadas por \texttt{psrecord}: el tiempo corresponde al último valor registrado en la columna \textit{Elapsed time}, mientras que la memoria máxima corresponde al mayor valor de la columna \textit{Real (MB)}. A partir de estas medidas se calcularon dos indicadores de lectura directa:

\[
\text{Ahorro de memoria (\%)} = 100\left(1-\frac{M_{\text{comp}}}{M_{\text{ref}}}\right),
\qquad
\text{Tiempo relativo (x)} = \frac{T_{\text{comp}}}{T_{\text{ref}}}.
\]

En estas expresiones, $M_{\text{ref}}$ y $T_{\text{ref}}$ corresponden a la ejecución con almacenamiento no comprimido del mismo circuito y del mismo número de qubits. Un valor negativo en el ahorro de memoria indica que la estrategia evaluada consumió más memoria que la referencia, lo que permite identificar con claridad los casos en los que la compresión no alcanza a compensar sus propios costos auxiliares.

\subsubsection{Verificación de exactitud numérica}

La exactitud numérica se evaluó según la naturaleza de cada estrategia. En Blosc, al tratarse de una compresión sin pérdida, la validación consistió en exportar el vector completo de amplitudes de cada ejecución comprimida y compararlo amplitud por amplitud con la ejecución no comprimida de la misma instancia, con el fin de verificar igualdad exacta respecto a la referencia. En ZFP, esa misma comparación completa se empleó para medir el máximo error absoluto, ya que en este caso el interés no era demostrar igualdad exacta, sino cuantificar la desviación introducida por la compresión con pérdida.

\subsubsection{Criterio de lectura de las figuras}

En todas las figuras, las líneas continuas representan resultados medidos hasta 25 qubits. A partir de ese punto, las líneas punteadas muestran proyecciones de tendencia hasta 28 qubits y las bandas sombreadas representan la incertidumbre asociada a esa estimación. Los fundamentos teóricos y los criterios estadísticos empleados para obtener estas predicciones se describen en el Apéndice~\ref{ap:proyeccion_graficas}. Estas proyecciones se utilizan únicamente como apoyo interpretativo para discutir la dirección esperada de las curvas si se ampliara el tamaño del problema dentro del mismo régimen experimental; por tanto, no sustituyen la evidencia medida ni se emplean como base de las conclusiones principales.

\subsubsection{Comparación complementaria para Grover}

Además de la comparación principal entre almacenamiento no comprimido, Blosc y ZFP, en Grover se incorporó una evaluación complementaria con la variante de parametrización reducida descrita en la Sección~\ref{subsec:grover_reducido_impl}. Esta comparación adicional permitió distinguir entre el beneficio derivado de comprimir el vector de estado denso y el beneficio asociado a una reformulación específica del algoritmo que reduce el volumen efectivo de datos manipulados. En esta subsección complementaria, la discusión se concentra en memoria y tiempo de ejecución, pues la exactitud ya queda determinada por la naturaleza exacta de la variante implementada y por la verificación de referencia efectuada en los experimentos de grover estándar.
